\chapter{RISC-V 启动与内存管理实现}\label{chap:riscv-mm}

启动与内存管理是操作系统移植中最基础的部分。对于原 IA-32 版本 Unix V6++ 而言，启动过程依赖传统 PC 的实模式、保护模式和引导扇区加载机制；而在 RISC-V/QEMU virt 平台上，内核通常由 QEMU 和 OpenSBI 共同引导，进入内核时已经处于 RISC-V 特权架构定义的运行环境中。因此，本课题需要重新设计内核入口、链接布局、早期栈、页表初始化和物理内存管理机制。本章围绕当前代码中的启动流程、Sv39 地址空间设计和页帧分配机制展开说明。

\section{内核启动流程}

重构后系统通过 QEMU 的 \verb|-kernel| 参数加载 RISC-V 内核镜像，并使用 OpenSBI 作为固件层。OpenSBI 负责完成机器模式下的基础初始化，并将控制权转交给内核。内核入口由链接脚本 \verb|src/kernel-riscv64.ld| 指定，入口符号为 \verb|_start|，加载地址为 \verb|0x80200000|。该地址位于 QEMU virt 平台常用的 DRAM 区域内，也避开了 OpenSBI 自身占用的低地址空间。

链接脚本将内核划分为 \verb|.bootstrap|、\verb|.text|、\verb|.rodata|、\verb|.data| 和 \verb|.bss| 等段。其中 \verb|.bootstrap.entry| 保存最早执行的入口代码，\verb|.bootstrap.stack| 保存启动阶段使用的静态栈空间，\verb|__bss_start|、\verb|__bss_end| 和 \verb|__kernel_end| 等符号则供 Rust 代码在运行时定位 BSS 段和内核结束位置。相比原 IA-32 启动代码中较长的实模式到保护模式切换流程，RISC-V 版本的早期启动逻辑更集中，主要任务是建立 Rust 运行所需的最小执行环境。

内核入口 \verb|_start| 使用 Rust 的裸函数和内联汇编实现。入口代码首先将栈指针 \verb|sp| 设置到预留的 bootstrap 栈顶，然后跳转到 \verb|riscv64_rust_entry|。OpenSBI 传入的 \verb|hart_id| 和设备树地址 \verb|dtb_addr| 保留在参数寄存器中，由 Rust 入口函数接收并输出到串口日志。当前系统以单核方式运行，QEMU 启动参数中使用 \verb|-smp 1|，因此启动流程只围绕一个 hart 展开。

\begin{table}[htbp]
  \centering
  \caption{RISC-V 内核早期启动阶段}\label{tab:riscv-boot-steps}
  \renewcommand{\arraystretch}{1.35}
  \begin{tabular}{p{3.2cm}p{10.2cm}}
    \toprule
    \textbf{阶段} & \textbf{主要工作} \\
    \midrule
    QEMU 加载内核 & 将内核 ELF 加载到 \verb|0x80200000|，并准备 virt 平台设备环境 \\
    OpenSBI 转交控制权 & 完成机器模式基础初始化，将 hart 编号和设备树地址传给内核 \\
    \verb|_start| 入口 & 设置 bootstrap 栈，跳转到 Rust 内核入口函数 \\
    BSS 清零 & 根据链接脚本符号清空 \verb|.bss|，保证静态数据初值正确 \\
    页表初始化 & 建立内核映射、用户页表入口和 MMIO 映射，写入 \verb|satp| 启用 Sv39 \\
    核心子系统初始化 & 初始化串口、页管理器、进程零号、Trap、缓冲区、设备和文件系统 \\
    进入调度器 & 创建初始进程并进入调度循环，后续由用户态 Shell 接管交互 \\
    \bottomrule
  \end{tabular}
\end{table}

\verb|riscv64_rust_entry| 中的初始化顺序反映了内核各模块之间的依赖关系。系统首先清零 BSS 段，随后初始化页目录和用户页表并启用页表保护；页表启用后初始化串口，使后续初始化过程具备可观察的日志输出；然后初始化物理页管理器、零号进程和 Trap 入口；再初始化缓冲区、外部中断控制器和定时器；最后加载文件系统超级块，设置根目录和控制终端，创建初始进程并进入调度器。这个流程将硬件平台初始化和 Unix V6++ 上层语义连接起来，是后续用户态程序能够运行的基础。

在 Rust 重写过程中，启动流程的一个直接收益是初始化步骤以函数调用形式清晰排列在入口函数中。汇编部分被限制在设置栈和跳转入口这一极小范围内，其余逻辑由 Rust 表达。这样既保留了裸机启动所需的底层控制能力，也使初始化顺序、错误处理和调试输出更容易阅读和维护。对于毕业设计实现而言，这部分工作不仅包括入口函数本身，还包括链接脚本、段符号、早期栈、串口日志和后续子系统初始化顺序的整体配合。

\section{地址空间与页表设计}

RISC-V 64 位平台常用 Sv39 分页模式。Sv39 使用三级页表，虚拟地址被划分为多级索引和页内偏移，页表项中包含有效位、读写执行权限、用户位、全局位、访问位和脏位等标志\cite{riscv_privileged_20211203}。当前系统在 \verb|src/machine/page_table.rs| 中定义了 \verb|PageDirectory|、\verb|PageDirectoryEntry|、\verb|PageTable| 和 \verb|PageTableEntry| 等结构，用 Rust 类型对页表项编码、权限标志和页表数组进行封装。

系统地址空间设计兼顾三类访问需求：第一，内核需要访问自身代码、数据和管理的物理内存；第二，用户程序需要拥有从低地址开始的用户空间；第三，内核需要访问 UART、PLIC、VirtIO 等 MMIO 设备。为此，当前页表同时建立低地址用户映射入口、物理内存恒等映射、内核高地址映射和 MMIO 映射。内核高地址从 \verb|0xc0000000| 开始，MMIO 区域从 \verb|0xf0000000| 开始，物理内存基址为 \verb|0x80000000|。

\begin{table}[htbp]
  \centering
  \caption{当前系统主要地址布局}\label{tab:address-layout}
  \renewcommand{\arraystretch}{1.35}
  \begin{tabular}{p{3.6cm}p{3.8cm}p{5.8cm}}
    \toprule
    \textbf{区域} & \textbf{地址或范围} & \textbf{用途} \\
    \midrule
    DRAM 物理基址 & \verb|0x80000000| & QEMU virt 平台内存起始地址 \\
    内核加载地址 & \verb|0x80200000| & RISC-V 内核 ELF 加载和链接起点 \\
    内核高地址基址 & \verb|0xc0000000| & 内核运行时访问物理内存的高地址窗口 \\
    内核 MMIO 基址 & \verb|0xf0000000| & UART、VirtIO、PLIC 等设备映射窗口 \\
    用户空间大小 & 约 8MB & 由 4 个用户页表覆盖，每个页表覆盖 2MB \\
    当前管理内存 & 64MB & 页表和页帧管理器当前纳入管理的物理内存范围 \\
    \bottomrule
  \end{tabular}
\end{table}

页表初始化由 \verb|init_page_directory| 和 \verb|init_user_page_table| 完成。前者清空根页表和二级页表，随后建立低地址页表入口、物理内存恒等映射和内核高地址映射；后者初始化用户页表数组，并将其挂接到低地址页表入口。系统随后调用 \verb|enable_page_protection|，将根页表物理页号写入 \verb|satp|，并执行 \verb|sfence.vma| 刷新 TLB，使新的地址翻译规则生效。

内核高地址映射采用大页映射与普通页表窗口结合的方式。对于大部分物理内存，系统以 2MB 为粒度建立直接映射，以减少页表项数量；对于需要细粒度切换的用户区窗口，则使用普通页表。当前代码中的 \verb|switch_user_struct| 会修改内核页表中固定位置的页表项，将当前进程的用户区物理页映射到内核可访问的窗口中，并刷新 TLB。这样可以保留 Unix V6++ 中类似 u-area 的访问方式，同时适配 RISC-V 的页表机制。

Rust 在页表实现中的作用主要体现在两个方面。其一，页表项权限通过 \verb|EntryFlags| 位标志表示，读、写、执行、用户、全局等权限组合比手写整数常量更直观。其二，\verb|PageTable| 实现了索引访问和迭代接口，页表初始化过程可以表达为对结构体数组的清空、挂接和映射，而不是散落在各处的地址计算。底层编码仍然需要精确遵守 RISC-V 页表项格式，但 Rust 类型封装降低了误用权限位和页表层级的风险。

\section{物理内存管理}

页表建立之后，内核需要管理可分配的物理页。当前系统在 \verb|src/mm/| 目录中实现物理页描述、内核页分配、用户页分配、内核栈和全局堆分配器等功能。物理页大小为 4KB，系统以 \verb|PhysPage| 描述每个页帧，并在 \verb|MemoryZone| 中维护页帧号与页描述结构之间的映射关系。当前实现将 64MB 物理内存纳入页帧管理，每个页帧对应一个 \verb|PhysPage| 结构。

\verb|init_page_managers| 是物理页管理的初始化入口。该函数通过链接脚本导出的 \verb|__kernel_end| 获得内核镜像结束地址，并按页对齐后作为可分配内核页区域的起点。系统将物理内存划分为内核页区域和用户页区域：内核页主要用于内核对象、内核堆、设备缓冲区和内核栈；用户页主要用于进程地址空间。两个区域分别由 \verb|KERNEL_PAGE_MANAGER| 和 \verb|USER_PAGE_MANAGER| 管理，底层均使用 Buddy 分配器。

\begin{table}[htbp]
  \centering
  \caption{物理内存管理相关模块}\label{tab:mm-modules}
  \renewcommand{\arraystretch}{1.35}
  \begin{tabular}{p{3.5cm}p{9.8cm}}
    \toprule
    \textbf{模块} & \textbf{职责} \\
    \midrule
    \verb|zone.rs| & 建立物理页描述数组，实现页帧号与 \verb|PhysPage| 的相互定位 \\
    \verb|page.rs| & 定义 \verb|PhysPage|、\verb|KernelPages|、\verb|UserPages| 和页列表适配器 \\
    \verb|page_manager.rs| & 初始化内核页和用户页 Buddy 分配器，提供页级分配与释放接口 \\
    \verb|allocator.rs| & 实现内核全局堆分配器，小对象使用 slab，大对象回退到页分配器 \\
    \verb|kstack.rs| & 为每个进程分配独立内核栈，并计算栈顶虚拟地址 \\
    \verb|swapper_manager.rs| & 维护交换区分配，用于继承 Unix V6++ 中换出相关结构 \\
    \bottomrule
  \end{tabular}
\end{table}

在页分配接口设计上，当前系统区分 \verb|KernelPages| 和 \verb|UserPages|。这两个类型共享底层 \verb|Pages| 泛型实现，但根据类型参数选择不同的分配器。更重要的是，\verb|Pages| 实现了 \verb|Drop|：当封装对象被释放时，其持有的物理页会自动归还给对应的 Buddy 分配器。这是 Rust 重写带来的一个实际好处。对于内核这种资源生命周期复杂的系统，手工释放容易遗漏，而将“页的所有权”和“页的释放动作”绑定到类型上，可以降低资源泄漏的概率。当然，进程地址空间、共享文本段和交换区等复杂场景仍然需要谨慎设计所有权转移，但类型封装已经为后续维护提供了更清晰的边界。

内核堆分配器由 \verb|allocator.rs| 实现，并通过 \verb|#[global_allocator]| 注册为 Rust 全局分配器。小于等于 2048 字节的分配请求由 slab 分配器处理，较大对象则从内核页分配器中申请连续页。这样，Rust 标准分配接口中的 \verb|Box|、\verb|Vec| 等容器可以在内核中使用，为进程表、文件对象、目录项缓存和其他动态结构提供基础。由于系统运行在 \verb|no_std| 环境下，这一全局分配器是 Rust 内核能够使用 \verb|alloc| crate 的前提之一。

进程内核栈也由内存管理模块统一分配。当前每个进程拥有独立的 \verb|KernelStack|，大小为 8 页，即 32KB。创建进程时，内核通过 \verb|KernelStack::new| 从内核页分配器申请栈空间，并通过物理地址到高地址虚拟地址的转换计算栈顶。独立内核栈配合 Trap 上下文和任务上下文切换，使每个进程在陷入内核时都有自己的内核执行空间，为后续进程调度和系统调用处理提供基础。

\section{本章小结}

本章介绍了 RISC-V/Rust 版本 Unix V6++ 的启动与内存管理实现。系统通过 QEMU/OpenSBI 进入位于 \verb|0x80200000| 的内核入口，使用少量汇编设置早期栈后进入 Rust 初始化函数；随后完成 BSS 清零、Sv39 页表建立、串口和核心子系统初始化，并进入进程调度。内存管理方面，系统设计了内核高地址映射、用户空间映射和 MMIO 映射，并使用 Buddy 分配器、slab 分配器和 Rust 全局分配器组织物理页与内核堆。

从工作量角度看，本章所述内容覆盖了跨平台移植中最底层的一组任务：重写 RISC-V 链接脚本和启动入口，建立可运行的 Rust 裸机环境，设计并实现 Sv39 页表结构，完成地址转换和 TLB 刷新接口，划分内核页与用户页区域，并将页分配器、内核堆和进程内核栈接入后续内核模块。这些工作为中断、系统调用、进程管理和文件系统运行奠定了基础。
